\documentclass[12pt, fleqn, openany]{studyGuide}
\setstretch{1.4}
\enableInlineReview
\renewcommand{\VMBookTitle}{Audit Review}
\renewcommand{\VMBookSubtitle}{Grade 7 --- Ch Changed Questions --- Changed Questions}
\renewcommand{\VMFooterURL}{https://viewmath.com/grade7}
\renewcommand{\VMFooterURLDisplay}{ViewMath.com/Grade7}
\begin{document}

\begin{center}
\begin{tcolorbox}[
	enhanced,
	colback=white,
	colframe=funTeal!60,
	boxrule=1.5pt,
	arc=10pt,
	outer arc=10pt,
	width=0.9\linewidth,
	halign=center,
	top=5mm, bottom=5mm,
]
	{\Large\bfseries\sffamily\color{funTealDark}Audit Review}\\[2mm]
	{\large\sffamily\color{funGrayDark}Grade 7 --- Ch Changed Questions --- Changed Questions}\\[2mm]
	{\normalsize\sffamily\color{funGrayDark}Verify corrections below}
\end{tcolorbox}
\end{center}

\newpage
\section*{8-3-q17 \hfill \normalsize\texttt{tests\_questions\_bank\_2/topics/ch08-03-comparing-two-populations-visually.tex}}
\begin{practiceQuestion}{8-3-q17}{sa}
\begin{questionText}
Team C: mean $65$, MAD $7$. Team D: mean $72$, MAD $7$. Is the difference in means meaningful? Justify using MADs.
\end{questionText}
\correctAnswer{No; the difference is $1$ MAD, so it is not clearly meaningful.}
\explanation{Difference: $72 - 65 = 7$. In MADs: $\frac{7}{7} = 1$. A difference of $1$ MAD shows only modest separation with considerable overlap, so it is not usually considered a meaningful difference.}
\end{practiceQuestion}

\newpage
\section*{8-3-q27 \hfill \normalsize\texttt{tests\_questions\_bank\_2/topics/ch08-03-comparing-two-populations-visually.tex}}
\begin{practiceQuestion}{8-3-q27}{gsa}
\begin{questionText}
The dot plots show hours of community service completed by members of two clubs.

\begin{center}
\begin{tikzpicture}[scale=0.5]
  % Club A number line
  \draw[thick,->] (3.5,0) -- (14,0);
  \foreach \x in {4,...,13} \draw (\x,0.15) -- (\x,-0.15) node[below,font=\tiny] {$\x$};
  \node[below] at (8.5,-0.9) {Hours};
  % Club A: 6,7,7,8,8,8,9,9,10 sum=72 mean=8
  \fill[funBlue] (6,0.5) circle (4pt);
  \foreach \y in {1,2} \fill[funBlue] (7,0.1+\y*0.4) circle (4pt);
  \foreach \y in {1,2,3} \fill[funBlue] (8,0.1+\y*0.4) circle (4pt);
  \foreach \y in {1,2} \fill[funBlue] (9,0.1+\y*0.4) circle (4pt);
  \fill[funBlue] (10,0.5) circle (4pt);
  \node[funBlue,above] at (8,2) {\small Club A};
  % Club B number line
  \draw[thick,->] (3.5,-3.3) -- (14,-3.3);
  \foreach \x in {4,...,13} \draw (\x,-3.15) -- (\x,-3.45) node[below,font=\tiny] {$\x$};
  % Club B: 9,10,10,11,11,11,12,12,13 sum=99 mean=11
  \fill[funRed] (9,-3.3+0.5) circle (4pt);
  \foreach \y in {1,2} \fill[funRed] (10,-3.3+0.1+\y*0.4) circle (4pt);
  \foreach \y in {1,2,3} \fill[funRed] (11,-3.3+0.1+\y*0.4) circle (4pt);
  \foreach \y in {1,2} \fill[funRed] (12,-3.3+0.1+\y*0.4) circle (4pt);
  \fill[funRed] (13,-3.3+0.5) circle (4pt);
  \node[funRed,below] at (11,-3.5) {\small Club B};
\end{tikzpicture}
\end{center}

Find the mean for each club. Then find the MAD for Club A and express the difference in means as a multiple of the MAD.
\end{questionText}
\correctAnswer{Club A mean $= 8$, Club B mean $= 11$; MAD of Club A $\approx 0.9$; difference $\approx 3.4$ MADs}
\explanation{Club A: $\frac{6+7+7+8+8+8+9+9+10}{9} = \frac{72}{9} = 8$. Deviations from $8$: $2,1,1,0,0,0,1,1,2$; MAD $= \frac{8}{9} \approx 0.89$. Club B mean: $\frac{9+10+10+11+11+11+12+12+13}{9} = \frac{99}{9} = 11$. Difference: $\frac{11-8}{0.89} \approx 3.4$ MADs, indicating a meaningful difference.}
\end{practiceQuestion}

\newpage
\section*{8-4-q12 \hfill \normalsize\texttt{tests\_questions\_bank\_2/topics/ch08-04-comparing-populations-with-measures.tex}}
\begin{practiceQuestion}{8-4-q12}{mc}
\begin{questionText}
Data set A: $20, 22, 24, 26, 28$. Data set B: $10, 18, 24, 30, 38$. Both sets have a mean of $24$. Which set has a larger IQR?
\end{questionText}
\choiceA{Set A}
\choiceB{Set B}
\choiceC{They have the same IQR}
\choiceD{IQR cannot be compared}
\correctAnswer{B}
\explanation{Set A: $Q_1 = 21$, $Q_3 = 27$, IQR $= 6$. Set B: $Q_1 = 14$, $Q_3 = 34$, IQR $= 20$. Set B has a much larger IQR, showing more variability.}
\end{practiceQuestion}

\newpage
\section*{8-4-q27 \hfill \normalsize\texttt{tests\_questions\_bank\_2/topics/ch08-04-comparing-populations-with-measures.tex}}
\begin{practiceQuestion}{8-4-q27}{gsa}
\begin{questionText}
The dot plots show the number of sit-ups completed by students in two fitness classes.

\begin{center}
\begin{tikzpicture}[scale=0.5]
  \draw[thick,->] (14.5,0) -- (30,0);
  \foreach \x in {15,...,29} \draw (\x,0.15) -- (\x,-0.15) node[below,font=\tiny] {$\x$};
  \node[below] at (22,-0.9) {Sit-ups};
  % Fitness A: 16,18,18,20,20,20,22,22,24 mean=180/9=20
  \fill[funBlue] (16,0.5) circle (4pt);
  \foreach \y in {1,2} \fill[funBlue] (18,0.1+\y*0.4) circle (4pt);
  \foreach \y in {1,2,3} \fill[funBlue] (20,0.1+\y*0.4) circle (4pt);
  \foreach \y in {1,2} \fill[funBlue] (22,0.1+\y*0.4) circle (4pt);
  \fill[funBlue] (24,0.5) circle (4pt);
  \node[funBlue,above] at (20,2) {\small Class A};
  % Fitness B: 20,22,22,24,24,24,26,26,28 mean=216/9=24
  \fill[funRed] (20,0.5) circle (4pt);
  \foreach \y in {1,2} \fill[funRed] (22,0.1+\y*0.4) circle (4pt);
  \foreach \y in {1,2,3} \fill[funRed] (24,0.1+\y*0.4) circle (4pt);
  \foreach \y in {1,2} \fill[funRed] (26,0.1+\y*0.4) circle (4pt);
  \fill[funRed] (28,0.5) circle (4pt);
  \node[funRed,above] at (24,2) {\small Class B};
\end{tikzpicture}
\end{center}

Find the mean for each class and the MAD for Class A. Then express the difference in means as a multiple of the MAD.
\end{questionText}
\correctAnswer{Class A mean $= 20$, Class B mean $= 24$; MAD of A $\approx 1.8$; difference $\approx 2.2$ MADs}
\explanation{Class A: $\frac{16+18+18+20+20+20+22+22+24}{9} = \frac{180}{9} = 20$. Deviations: $4,2,2,0,0,0,2,2,4$; MAD $= \frac{16}{9} \approx 1.78$. Class B: $\frac{20+22+22+24+24+24+26+26+28}{9} = \frac{216}{9} = 24$. Difference: $\frac{24-20}{1.78} \approx 2.25$ MADs, suggesting a meaningful difference.}
\end{practiceQuestion}

\newpage
\section*{9-2-q24 \hfill \normalsize\texttt{tests\_questions\_bank\_2/topics/ch09-02-theoretical-probability.tex}}
\begin{practiceQuestion}{9-2-q24}{gmc}
\begin{questionText}
A dart is thrown randomly at the target below. The target has $4$ concentric rings of equal area, colored from inside to outside.

\begin{center}
\begin{tikzpicture}
  \fill[funOrange!50] (0,0) circle (4);
  \fill[funGreen!50] (0,0) circle (3.4641);
  \fill[funBlue!50] (0,0) circle (2.8284);
  \fill[funRed!50] (0,0) circle (2);
  \node at (0,0) {Red};
  \node at (0,2.35) {Blue};
  \node at (0,3.15) {Green};
  \node at (0,3.72) {Orange};
\end{tikzpicture}
\end{center}

If each ring has equal area, what is the probability that the dart lands in the green ring?
\end{questionText}
\choiceA{$\dfrac{1}{2}$}
\choiceB{$\dfrac{1}{4}$}
\choiceC{$\dfrac{1}{3}$}
\choiceD{$\dfrac{3}{4}$}
\correctAnswer{B}
\explanation{Since each of the $4$ rings has equal area and the dart lands randomly, this is a uniform model. The probability of landing in any one ring is $\frac{1}{4}$.}
\end{practiceQuestion}

\newpage
\section*{3-9-q17 \hfill \normalsize\texttt{tests\_questions\_bank\_2/topics\_additional/ch03-09-introduction-to-square-roots.tex}}
\begin{practiceQuestion}{3-9-q17}{sa}
\begin{questionText}
Estimate $\sqrt{72}$ to the nearest whole number.
\end{questionText}
\correctAnswer{$8$}
\explanation{$8^2 = 64$ and $9^2 = 81$. Since $72 - 64 = 8$ and $81 - 72 = 9$, $72$ is closer to $64$, so $\sqrt{72} \approx 8.49$ and rounds to $8$.}
\end{practiceQuestion}

\newpage
\section*{4-7-q22 \hfill \normalsize\texttt{tests\_questions\_bank\_2/topics\_additional/ch04-07-laws-of-exponents.tex}}
\begin{practiceQuestion}{4-7-q22}{gmc}
\begin{questionText}
The diagram below shows an exponent rule machine. What is the output?

\begin{center}
\begin{tikzpicture}[scale=0.7]
  \node[draw, thick, rounded corners, fill=funBlue!15, minimum width=2.5cm, minimum height=1cm] (input) at (0,0) {Input: $6^5$};
  \node[draw, thick, rounded corners, fill=funOrange!20, minimum width=3cm, minimum height=1.5cm, align=center] (machine) at (5,0) {MULTIPLY \\ by $6^3$};
  \node[draw, thick, rounded corners, fill=funGreen!15, minimum width=2.5cm, minimum height=1cm] (output) at (10,0) {Output: $?$};
  \draw[->, thick] (input) -- (machine);
  \draw[->, thick] (machine) -- (output);
\end{tikzpicture}
\end{center}
\end{questionText}
\choiceA{$6^{15}$}
\choiceB{$6^8$}
\choiceC{$6^2$}
\choiceD{$36^8$}
\correctAnswer{B}
\explanation{Product rule: $6^5 \times 6^3 = 6^{5+3} = 6^8$.}
\end{practiceQuestion}

\newpage
\section*{8-5-q18 \hfill \normalsize\texttt{tests\_questions\_bank\_2/topics\_additional/ch08-05-stem-and-leaf-plots.tex}}
\begin{practiceQuestion}{8-5-q18}{sa}
\begin{questionText}
Create a stem-and-leaf plot (list the stems and their leaves) for this data: $25, 31, 28, 34, 22, 30, 27, 35, 29$.
\end{questionText}
\correctAnswer{Stem $2$: $2, 5, 7, 8, 9$; Stem $3$: $0, 1, 4, 5$}
\explanation{Sort the data: $22, 25, 27, 28, 29, 30, 31, 34, 35$. Stems are the tens digits. Stem $2$: leaves $2, 5, 7, 8, 9$. Stem $3$: leaves $0, 1, 4, 5$ (listed in order).}
\end{practiceQuestion}

\newpage
\section*{m-6-6-q25 \hfill \normalsize\texttt{tests\_questions\_bank\_2/topics\_modified/ch06-06-angle-relationships.tex}}
\begin{practiceQuestion}{m-6-6-q25}{gsa}
\begin{questionText}
In the diagram, two parallel lines $\ell_1$ and $\ell_2$ are cut by a transversal $t$. Find the values of $a$ and $b$.

\begin{center}
\begin{tikzpicture}[scale=0.8]
  % Parallel lines
  \draw[thick] (-3,2) -- (4,2) node[right] {$\ell_1$};
  \draw[thick] (-3,-1) -- (4,-1) node[right] {$\ell_2$};
  % Transversal
  \draw[thick] (-1,3.5) -- (2,-2.5) node[right] {$t$};
  % Angle labels
  \node at (1.2,2.4) {$a°$};
  \node at (-0.2,1.5) {$115°$};
  \node at (1.8,-0.5) {$b°$};
\end{tikzpicture}
\end{center}
\end{questionText}
\correctAnswer{$a = 115°$, $b = 115°$}
\explanation{The $115°$ angle and angle $a$ are vertical angles at the top intersection, so $a = 115°$. Angle $b$ is corresponding to angle $a$ on the lower parallel line, so $b = 115°$.}
\end{practiceQuestion}

\end{document}